\section{Two informative raw-kernel realizations}\label{sec:realizations}
Both constructions start with a compact $d$-regular probability $\nu_K$ and a reliability law $\rho_\alpha(dt)=\alpha t^{\alpha-1}dt$ on $(0,1]$. The measure may have the singular location support just described. Their actual prepared reference reports have law
\[
 \nu(dy,dt,dv)=\nu_K(dy)\rho_\alpha(dt)\tfrac12(\delta_{-1}+\delta_1)(dv).
\]
The common proxy and risk mark are $f(y,t,v)=y$ and $w(y,t,v)=t^2$. The protected task is vector prediction, not parameter estimation.

\begin{theorem}[A controlled hidden-state sensor]\label{thm:classical}
Let a hidden sign $E$ evolve, before each report, by one of finitely many controlled transition matrices. Every entry lies in $[q,1-q]$, where $0<q\le1/2$; the matrices may belong to an unknown fixed family. Independently draw $(Y,T)$ with law $\nu_K\otimes\rho_\alpha$, and report $V\in\{-1,1\}$ with
\[
 \Pp(V=v\mid E=e)=\frac{1+\eta ve}{2},\qquad0<\eta<1.
\]
The current physical state includes $(E,Y,T)$. After stopping, audit
$Z=Y+TEe_1$, where $e_1$ is a fixed unit vector. No transition occurs before this audit. The observer sees $(Y,T,V)$, not $E$.

Theorem~\ref{thm:central} applies uniformly to the whole transition family, and the matching law \eqref{eq:crossover-aut} holds. A reference member with a symmetric transition and initial fair sign gives the stated actual first-call law. Rejected calls do not reset the sign, and the conditional audit mean is generally not a function of the last report alone.
\end{theorem}
\begin{proof}
Set $\gamma=(1-\eta)/2$ and $a=q\gamma/(1-\gamma)>0$. Before each report, the conditional probability of either sign is at least $q$, whatever the surviving history and chosen transition. Bayes' rule after $V$ gives probabilities in $[a,1-a]$. The report likelihood relative to the fair $V$ reference is between $2\gamma$ and $2(1-\gamma)$, and $Y,T$ retain their displayed conditional law. Therefore \eqref{eq:report-envelope} holds at every history with $\ell=2\gamma$, $u=2(1-\gamma)$.

Writing $b=\E[E\mid\mathcal H_{t-1},Y,T,V]$, the audit mean is $Y+Tb e_1$ and its variance is $T^2(1-b^2)$. Hence
\[
 4a(1-a)T^2\le v\le T^2,\qquad \|m-Y\|^2\le T^2.
\]
These prove the raw calibration certificate with $B_0=0$, $c_v=4a(1-a)$, $C_v=c_b=1$. They can equivalently be checked by multiplying through the positive report likelihood as in \eqref{eq:raw-certificate}. The marked law is the actual law of $(Y,T)$, so \eqref{eq:rectangle-mass} follows from $d$-regularity and the exact mass $s^\alpha$.

For completeness, the predictive carrier is obtained from the raw two-coordinate positive instrument, with the reported $Y,T$ included for a possible immediate audit. Future instruments act linearly on the unnormalized two sign masses; audit probabilities and their moments are executable functions of that state. Distinct sign beliefs are separated by a future $V$ report when the chosen transition is invertible; when it is not, the redundant direction is quotiented out rather than declared identifiable. An invertible transition exists in the nontrivial reference subclass $q<1/2$. The last reported $Y,T$ provide a Borel realization of the immediate-task coordinates; the full predictive quotient is not the proxy $Y$.

Report total variation is continuous, indeed Lipschitz, in the incoming sign belief. Bayesian updates are ratios with denominators at least $\gamma$, so are uniformly continuous on the compact belief interval. Positive transitions supply an invariant interval before each report. The risk proof uses this invariant calibration, not a stored posterior recursion or an uncharged numerical register. It remains valid for any causal transition selection respecting the same raw entry bounds.
\end{proof}

\begin{theorem}[Noncommuting weak quantum sensors]\label{thm:quantum}
Prepare a qubit in $I/2$. Let
\[
 D_q(\rho)=(1-2q)\rho+q\,\tr(\rho)I,
\qquad0<q\le1/2.
\]
For action $a$ choose a known unit vector $n_a\in\R^3$, a unitary $U_a$, and $0<\eta<1$. Write $\sigma=(\sigma_x,\sigma_y,\sigma_z)$ for the Pauli matrices and define
\[
 K_{a,v}=\left[\frac{I+\eta v\,n_a\cdot\sigma}{2}\right]^{1/2},
 \qquad
 \mathcal I_{a,v}(\rho)=K_{a,v}U_aD_q(\rho)U_a^*K_{a,v}.
\]
Independently report $(Y,T)$ with the same law as above. Retain the normalized post-instrument physical state. After commitment, measure $\sigma_z$ on that state to obtain $B\in\{-1,1\}$ and audit $Z=Y+TB e_1$.

These raw instruments verify the same Theorem~\ref{thm:central} and the same crossover \eqref{eq:crossover-aut}. The directions $n_a$ may be nonparallel and the instruments need not commute. Uniform bounds permit unknown $q$, directions or channels in a fixed admissible family; the implementing proxy program does not read the true member.
\end{theorem}
\begin{proof}
Each operation is completely positive and $\sum_vK_{a,v}^2=I$, so the instrument is normalized. On a density input, $D_q(\rho)\succeq qI$. With $\gamma=(1-\eta)/2$, the effect eigenvalues lie in $[\gamma,1-\gamma]$. Thus every report probability lies in this interval, and every normalized output state satisfies
\[
 \rho'\succeq \frac{q\gamma}{1-\gamma}I=aI.
\]
The protected Pauli outcome has variance $1-\tr(\sigma_z\rho')^2\ge4a(1-a)$. Its mean has absolute value at most one. Multiplication by $T$ gives exactly the same calibration constants as in Theorem~\ref{thm:classical}. From the specified preparation, the first $V$ is fair, because $D_q(I/2)=I/2$ and $\tr K_{a,v}^2=1$. Thus the reference report law is physically obtained, not assigned to a formal orbit.

Positive instrument operations on Hermitian matrices, and Born audit effects, give an unnormalized finite linear representation directly. Include a tomographically complete finite set of probe directions among executable future tests; their probabilities separate density operators on the attained image when $q<1/2$. A Borel realization is the compact density-operator body with the visible $Y,T$ audit registers; a redundant direction at a degenerate parameter is factored out. Report total variation is bounded by half the input trace distance after the channel. Normalization has denominator at least $\gamma$, so predictive update and report continuity are explicit on the invariant compact state body. No infinite-dimensional operator closure is used.

If two directions are nonparallel, the associated effects generally fail to commute. For instance, choosing $n_a$ on the $x$ and $z$ axes gives effects with nonzero commutator. Coherence can persist because $q<1/2$ and the effects are invertible; the operation is not a measure-and-prepare reset. The external observer nonetheless has only the charged classical labels. The physical qubit is not available as free observer memory. Finally the acquired $(Y,T)$ law proves \eqref{eq:rectangle-mass}, independently of the classical proof.
\end{proof}

The reliability mark is a measurable experimental quantity, not a promise that the decoder knows the hidden residual. In both models, replacing all past histories by the last visible proxy generally changes the posterior. What survives is the two-sided conditional risk comparison, uniformly after every rejected call. That is the hypothesis verified here and used by the general theorem.
